% =====================================================================
%  Capítulo 15 — Conclusiones
% =====================================================================
\chapter{Conclusiones}
\citacap{La mejor protección que puede tener una mujer es el valor.}{Elizabeth Cady Stanton}

Llegadas a este punto, ya deberíamos tener una idea de:

\begin{itemize}
  \item En qué circunstancias es oportuno defenderse de otra persona y cuál es el marco legal.
  \item Nuestras propias capacidades para ejecutar unas técnicas, descartar otras o entrenarlas más a fondo.
  \item Mediante la práctica, haber seleccionado movimientos preferidos para cada distancia de conflicto.
  \item Que los ataques deben efectuarse en combinación de técnicas, en series encadenadas y, a ser posible, a diferentes alturas.
  \item Que es de suma importancia la confianza en una misma: la duda genera un ataque débil que probablemente pondrá más violento al agresor. Tomada una decisión, no echarse atrás.
  \item Que se van a experimentar sensaciones corporales poco usuales, y que en situación real se realizan menos golpes, pero más contundentes.
  \item Que debemos estar preparadas mentalmente para recibir golpes o sufrir heridas, sabiendo que el organismo estará en un estado de mayor rendimiento y resistencia al dolor.
  \item Que el torrente hormonal del enfrentamiento es el cuerpo en su máximo punto de rendimiento; y que si alguna vez aparece la parálisis, es una respuesta automática que el entrenamiento ayuda a superar, nunca una culpa.
\end{itemize}

La defensa personal puede perfeccionarse día a día con la observación, la práctica física y el estudio. Lo más importante es la prevención y, llegado el enfrentamiento, la sorpresa (fingirse colaborativa o temerosa y no exteriorizar el ataque que se está por llevar a cabo). Y por último: no merece la pena sufrir daños por defender un bien material.

\section{Teléfonos y recursos esenciales (España)}

\begin{table}[h]
\centering
\renewcommand{\arraystretch}{1.4}
\begin{tabularx}{\linewidth}{@{}>{\bfseries\color{primario}}l X@{}}
\toprule
\textbf{112} & Emergencias generales (UE) \\
\textbf{091} & Policía Nacional \\
\textbf{062} & Guardia Civil \\
\textbf{016 / WhatsApp 600 000 016} & Violencia contra la mujer: gratuito, 24\,h, confidencial, no aparece en factura \\
\textbf{AlertCops (app)} & Chat y botón SOS con Policía y Guardia Civil \\
\textbf{Oficinas de Asistencia a las Víctimas} & Apoyo jurídico y psicológico gratuito (Ministerio de Justicia) \\
\bottomrule
\end{tabularx}
\end{table}
